\clearpage\fancyhead[L]{\small NEUE PALLAS-ELEMENTE.}\SourceLeaf{325}

\begin{center}\textbf{NEUE PALLAS-ELEMENTE.}\par\vspace{4pt}\hrule\vspace{4pt}\fontsize{10}{12}\selectfont Monatliche Correspondenz. Band XXII. S. 400\,..\,403. 1810 Oktober.\par\vspace{4pt}\hrule\vspace{4pt}\end{center}

\par\indent Von den neuen Planeten wird \emph{Pallas} zuerst sichtbar sein und wegen ihrer großen Erdnähe eine beträchtliche Helligkeit erreichen. In unseren nördlichen Breiten könnte diese allenfalls durch die starke südliche Deklination gemindert werden. Deshalb beeilen wir uns, unseren Lesern die Ephemeride ihrer Bewegung mitzuteilen, die Professor \textsc{Gauss} kürzlich berechnet hat. Anfang November wird sich der Planet wahrscheinlich bereits im Meridian beobachten lassen.

Die Elemente, auf denen die folgende Ephemeride beruht, wurden aus den Oppositionen von 1805, 7, 8 und 9 hergeleitet und sind nachstehend aufgeführt.

\Needspace{14\baselineskip}\begin{center}\emph{Göttinger Meridian.}\par\medskip\fontsize{10}{12}\selectfont\setlength{\tabcolsep}{4pt}\renewcommand{\arraystretch}{1.15}\begin{tabular}{llr}
Mittlere Länge 1803 & .\;.\;. & $221^\circ\;23'\;24''6$\\
mittlere tägliche Bewegung &  & $770''9365$\\
1803 Perihelium & .\;.\;.\;. & $120^\circ\;58'\;4''8$\\
--- $\AscendingNode$ & .\;.\;.\;.\;.\;. & $172\;27\;52.4$\\
Neigung & .\;.\;,\;.\;.\;.\;. & $34^\circ\;36'\;49''4$\\
Exzentrizität & .\;.\;.\;.\;. & $0.2446335$\\
Log. halbe Achse & .\;.\;.\;. & $0.4420473$\\
\end{tabular}\end{center}

\par\indent [Der Lauf der Pallas 1810 und 1811 ist für Mitternacht in Göttingen in Abständen von vier Tagen angegeben, Rektasz. und Dekl. in Minuten Log. der Entf. von der Erde mit vier Dezimalstellen]
